\documentclass[12pt,aspectratio=1610,professionalfont,serif,t]{beamer}
\usepackage[osf,p]{newpxtext}
\usepackage{newpxmath}
\usepackage{fontawesome5}

\usepackage{fancyvrb}
\fvset{%frame=single,
commandchars=\\\{\},
%framesep=1mm,
fontfamily=helvetica,
fontsize=\normalsize
}
%\usepackage[bitstream-charter]{mathdesign}
\usepackage{listings}
\lstset{commentstyle=\color{orange}, keywordstyle=\color{yellow} }


\usetheme{Boadilla}
\useoutertheme{split}
\usecolortheme{albatross}


\setbeamertemplate{blocks}[rounded][shadow=false]
\setbeamersize{text margin left=\dimexpr 1cm+1ex\relax, text margin right=\dimexpr 1cm+1ex\relax}
\addtobeamertemplate{headline}{}{\usebeamerfont{frametitle}\vspace*{1ex}}
\setbeamertemplate{navigation symbols}{}
\setbeamertemplate{page number in head/foot}{}
%\beamerdefaultoverlayspecification{<+->}


\setbeamercolor*{structure}{fg=green!60!black,bg=black}
\setbeamercolor*{normal text}{fg=green!60!black,bg=black}
\setbeamercolor{palette primary}{use={structure,normal text},fg=green,bg=structure.bg}
\setbeamercolor{palette secondary}{use={structure,normal text},fg=structure.fg,bg=structure.bg}
\setbeamercolor{palette tertiary}{use={structure,normal text},fg=structure.fg,bg=structure.bg}
\setbeamercolor{palette quaternary}{use={structure,normal text},fg=white,bg=green!30!black}
\setbeamercolor{palette sidebar quaternary}{use={structure,normal text},fg=white,bg=green!30!black}
\setbeamercolor{title in head/foot}{use={structure,normal text},fg=white,bg=green!30!black}

\setbeamercolor*{example text}{use={structure},fg=green!60!black}
\setbeamercolor*{block body}{use={structure},bg=structure.bg}
\setbeamercolor*{block body alerted}{use={structure},bg=structure.bg}
\setbeamercolor*{block body example}{use={structure},bg=structure.bg}
\setbeamercolor*{block title}{parent=structure,bg=structure.bg}
\setbeamercolor*{block title alerted}{use={structure,alerted text},fg=alerted text.fg,bg=structure.bg}
\setbeamertemplate{navigation symbols}{}
\setbeamertemplate{items}[square]
\setbeamercolor{item projected}{fg=white}
\setbeamercolor*{normal text}{use={structure},fg=white!90!blue,bg=structure.bg}
\setbeamercolor*{separation line}{}
\setbeamercolor*{fine separation line}{}
\setbeamercolor{alerted text}{use={structure},fg=structure.fg}

\usepackage[italian]{babel}
\usepackage[utf8]{inputenc}
\usepackage{pgf}
\usepackage{verbatim}
\usepackage{inconsolata}
\usepackage{listings}
\lstset{language=C, frame=single,
	basicstyle=\ttfamily,
	numbers=left, numberstyle=\tiny\color{gray},
	numbersep=5pt, fancyvrb=true
}

\usepackage[noend]{algorithm2e}
\usepackage{pgf}
\usepackage{tikz}
\usepackage{graphicx}
\usetikzlibrary{%
	arrows,
	arrows.meta,
	positioning,
	calc,
	backgrounds,
	chains,
	matrix,
	patterns,
	automata,
	fit,
	graphs,
	decorations,
	decorations.pathmorphing,
	decorations.pathreplacing,
	decorations.markings,
}

\usepgflibrary{shapes,shapes.geometric}


\usepackage{url}


\usepackage{xmpmulti}
\usepackage{transparent}
\graphicspath{{img/}}

\newcommand{\zero}{\ensuremath{\mathtt{0}}\xspace}
\newcommand{\uno}{\ensuremath{\mathtt{1}}\xspace}
\newcommand{\due}{\ensuremath{\mathtt{2}}\xspace}
\newcommand{\tre}{\ensuremath{\mathtt{3}}\xspace}
\newcommand{\kleene}{\ensuremath{{}^{\ast}}\xspace}
\xspaceaddexceptions{\zero \uno \due \tre \kleene}


\newcommand{\simplebox}[1]{
      \begin{tikzpicture}
              \node[text centered, rectangle, draw=black, thick, text width=\linewidth,align=center, rounded corners, minimum height=1em] (C) {#1};
      \end{tikzpicture}
}
%
% In questo file si mettono i comandi utili solo per questo insegnamento.
%
\author{Gianluca Della Vedova}
\title{Linguaggi e Computabilità}
\institute{Univ. Milano--Bicocca\\
	\texttt{https://www.unimib.it/gianluca-della-vedova}}
\date{\today}



\tikzset{
        albero/.style={level distance=9mm, every node/.style={font=\small},
                edge from parent/.style={draw, white!70!black}},
}

\begin{document}


\begin{frame}[fragile]
        \frametitle{Linguaggi e Computabilità: Lezione 16}
\begin{itemize}
\item
Grammatiche ambigue
\end{itemize}
\end{frame}




\begin{frame}
        \frametitle{Esempio 1: espressioni aritmetiche}
        \[
                E \to E + E \mid E \ast E \mid ( E ) \mid \mathtt{x}
        \]
        La parola $\mathtt{x} + \mathtt{x} \ast \mathtt{x}$ ha due alberi di derivazione:

        \medskip
        \begin{columns}[T]
                \begin{column}{0.45\textwidth}
                        \centering
                        \begin{tikzpicture}[albero, sibling distance=12mm]
                                \node {$E$}
                                child {node {$E$} child {node {$\mathtt{x}$}}}
                                child {node {$+$}}
                                child {node {$E$}
                                        child {node {$E$} child {node {$\mathtt{x}$}}}
                                        child {node {$\ast$}}
                                        child {node {$E$} child {node {$\mathtt{x}$}}}};
                        \end{tikzpicture}

                        $\mathtt{x} + (\mathtt{x} \ast \mathtt{x})$
                \end{column}
                \begin{column}{0.45\textwidth}
                        \centering
                        \begin{tikzpicture}[albero, sibling distance=12mm]
                                \node {$E$}
                                child {node {$E$}
                                        child {node {$E$} child {node {$\mathtt{x}$}}}
                                        child {node {$+$}}
                                        child {node {$E$} child {node {$\mathtt{x}$}}}}
                                child {node {$\ast$}}
                                child {node {$E$} child {node {$\mathtt{x}$}}};
                        \end{tikzpicture}

                        $(\mathtt{x} + \mathtt{x}) \ast \mathtt{x}$
                \end{column}
        \end{columns}
\end{frame}


% \begin{frame}
%         \frametitle{Esempio 1: derivazioni sinistre}
%         Prima derivazione (si applica prima $E \to E + E$):
%         \[
%                 E \Rightarrow E + E \Rightarrow \mathtt{x} + E
%                 \Rightarrow \mathtt{x} + E \ast E
%                 \Rightarrow \mathtt{x} + \mathtt{x} \ast E
%                 \Rightarrow \mathtt{x} + \mathtt{x} \ast \mathtt{x}
%         \]
%         Seconda derivazione (si applica prima $E \to E \ast E$):
%         \[
%                 E \Rightarrow E \ast E \Rightarrow E + E \ast E
%                 \Rightarrow \mathtt{x} + E \ast E
%                 \Rightarrow \mathtt{x} + \mathtt{x} \ast E
%                 \Rightarrow \mathtt{x} + \mathtt{x} \ast \mathtt{x}
%         \]

%         \bigskip
%         Grammatica non ambigua per lo stesso linguaggio (precedenza e associatività):
%         \[
%                 E \to E + T \mid T \qquad
%                 T \to T \ast F \mid F \qquad
%                 F \to ( E ) \mid \mathtt{x}
%         \]
% \end{frame}


\begin{frame}
        \frametitle{Esempio 2: \emph{dangling else}}
        \[
                S \to \texttt{if}\ b\ \texttt{then}\ S
                \mid \texttt{if}\ b\ \texttt{then}\ S\ \texttt{else}\ S
                \mid a
        \]
        La parola
        \[
                \texttt{if}\ b\ \texttt{then}\ \texttt{if}\ b\ \texttt{then}\ a\ \texttt{else}\ a
        \]
        ha due alberi di derivazione:
        \begin{itemize}
                \item
                $\texttt{if}\ b\ \texttt{then}\ \bigl(\texttt{if}\ b\ \texttt{then}\ a\ \texttt{else}\ a\bigr)$:
                l'\texttt{else} appartiene all'\texttt{if} interno
                \item
                $\texttt{if}\ b\ \texttt{then}\ \bigl(\texttt{if}\ b\ \texttt{then}\ a\bigr)\ \texttt{else}\ a$:
                l'\texttt{else} appartiene all'\texttt{if} esterno
        \end{itemize}
\end{frame}


\begin{frame}
        \frametitle{Esempio 2: alberi di derivazione}
        \begin{columns}[T]
                \begin{column}{0.5\textwidth}
                        \centering
                        \begin{tikzpicture}[albero, level 1/.style={sibling distance=11mm},
                                level 2/.style={sibling distance=8mm}]
                                \node {$S$}
                                child {node {\texttt{if}}}
                                child {node {$b$}}
                                child {node {\texttt{then}}}
                                child {node {$S$}
                                        child {node {\texttt{if}}}
                                        child {node {$b$}}
                                        child {node {\texttt{then}}}
                                        child {node {$S$} child {node {$a$}}}
                                        child {node {\texttt{else}}}
                                        child {node {$S$} child {node {$a$}}}};
                        \end{tikzpicture}
                \end{column}
                \begin{column}{0.5\textwidth}
                        \centering
                        \begin{tikzpicture}[albero, level 1/.style={sibling distance=11mm},
                                level 2/.style={sibling distance=8mm}]
                                \node {$S$}
                                child {node {\texttt{if}}}
                                child {node {$b$}}
                                child {node {\texttt{then}}}
                                child {node {$S$}
                                        child {node {\texttt{if}}}
                                        child {node {$b$}}
                                        child {node {\texttt{then}}}
                                        child {node {$S$} child {node {$a$}}}}
                                child {node {\texttt{else}}}
                                child {node {$S$} child {node {$a$}}};
                        \end{tikzpicture}
                \end{column}
        \end{columns}
\end{frame}


\begin{frame}
        \frametitle{Esempio 3: parentesi bilanciate}
        \[
                S \to S S \mid ( S ) \mid \varepsilon
        \]
        La parola $()()()$ ha (almeno) due alberi di derivazione:

        \medskip
        \begin{columns}[T]
                \begin{column}{0.45\textwidth}
                        \centering
                        \begin{tikzpicture}[albero, level distance=8mm, level 1/.style={sibling distance=30mm},
                                level 2/.style={sibling distance=16mm}, level 3/.style={sibling distance=6mm}]
                                \node {$S$}
                                child {node {$S$}
                                        child {node {$S$} child {node {$($}} child {node {$S$} child {node {$\varepsilon$}}} child {node {$)$}}}
                                        child {node {$S$} child {node {$($}} child {node {$S$} child {node {$\varepsilon$}}} child {node {$)$}}}}
                                child[level 2/.style={sibling distance=6mm}] {node {$S$} child {node {$($}} child {node {$S$} child {node {$\varepsilon$}}} child {node {$)$}}};
                        \end{tikzpicture}
                \end{column}
                \begin{column}{0.45\textwidth}
                        \centering
                        \begin{tikzpicture}[albero, level distance=8mm, level 1/.style={sibling distance=30mm},
                                level 2/.style={sibling distance=16mm}, level 3/.style={sibling distance=6mm}]
                                \node {$S$}
                                child[level 2/.style={sibling distance=6mm}] {node {$S$} child {node {$($}} child {node {$S$} child {node {$\varepsilon$}}} child {node {$)$}}}
                                child {node {$S$}
                                        child {node {$S$} child {node {$($}} child {node {$S$} child {node {$\varepsilon$}}} child {node {$)$}}}
                                        child {node {$S$} child {node {$($}} child {node {$S$} child {node {$\varepsilon$}}} child {node {$)$}}}};
                        \end{tikzpicture}
                \end{column}
        \end{columns}

        \medskip
        Peggio: anche $\varepsilon$ ha infiniti alberi
        ($S \Rightarrow SS \Rightarrow S \Rightarrow \varepsilon$, \dots).
\end{frame}


\begin{frame}
        \frametitle{Esercizio 1}
        Scrivere una grammatica context-free che genera il linguaggio
        \[
                L = \{ \zero^{n} \uno^{m} \; \vert \; n \neq m, n\ge 0, m\ge 0 \}
        \]
        sull'alfabeto $\Sigma_{B}$.
% Soluzione:
% S -> 0 S 1 | A | B
% A -> 0 A | 0
% B -> 1 B | 1
% S genera coppie 0...1 bilanciate, poi un blocco non vuoto di soli 0
% (n > m) oppure di soli 1 (n < m).
% La grammatica non è ambigua.
\end{frame}


\begin{frame}
        \frametitle{Esercizio 2}
        Scrivere una grammatica context-free che genera il linguaggio
        \[
                L = \{ \zero^{i} \uno^{j} \due^{k} \; \vert \; j = i + k, i>0, j>0 \}
        \]
        sull'alfabeto $\{\zero, \uno, \due\}$.
% Soluzione:
% S -> A B
% A -> 0 A 1 | epsilon
% B -> 1 B 2 | epsilon
% Una parola di L si scompone in modo unico come 0^i 1^i 1^k 2^k:
% A genera la prima parte, B la seconda.
\end{frame}


\begin{frame}
        \frametitle{Esercizio 3}
        Scrivere una grammatica context-free ambigua e una non ambigua che generino il linguaggio
        \[
                L = \{ \zero^{n} \uno^{m} \; \vert \; n \le m \le 2n, n>0, m>0 \}
        \]
        sull'alfabeto $\Sigma_{B}$.

        \bigskip

% Soluzione (ambigua):
% S -> 0 S 1 | 0 S 1 1 | 01 | 011
% Ogni 0 è accoppiato con uno oppure due 1.
% È ambigua: 00111 si ottiene come 0(0 1 1)1 oppure 0(0 1)11, cioè
% S => 0S1 => 00S111 => 00111  e  S => 0S11 => 00S111 => 00111.
% Soluzione (non ambigua): si usano prima tutte le coppie 0-1, poi le 0-11
% S -> 0 S 1 | T | 01
% T -> 0 T 1 1 |  011
\end{frame}


\end{document}

%%% Local Variables:
%%% mode: LaTeX
%%% TeX-engine: luatex
%%% TeX-master: t
%%% End:
